% ==============================================================================
% chapters/ch01_introduction.tex — บทที่ 1
% ==============================================================================
\chapter{บทนำและสถาปัตยกรรมระบบเกษตรอัจฉริยะ}
\label{ch:introduction}

ความมั่นคงทางอาหารและประสิทธิภาพการผลิตทางการเกษตรในคริสต์ศตวรรษที่ 21 กำลังเผชิญกับความท้าทายรอบด้าน ทั้งจากการเปลี่ยนแปลงสภาพภูมิอากาศ ความเสื่อมโทรมของทรัพยากรดิน และการระบาดของโรคและแมลงศัตรูพืช การเกษตรแม่นยำ (Precision Agriculture) และฟิสิกส์เกษตรดิจิทัล (Digital Agriphysics) จึงกลายเป็นระเบียบวิธีวิจัยและแนวทางปฏิบัติการสำคัญที่นำเทคโนโลยีเซนเซอร์ โทรมาตร (Telemetry) และปัญญาประดิษฐ์ (Artificial Intelligence) มาบูรณาการร่วมกันเพื่อเพิ่มผลผลิตและลดการใช้สารเคมีอย่างมีนัยสำคัญ

สถาปัตยกรรมระบบ \textbf{Embedded AI \& Edge AIoT for Agriculture and Environment} ถูกออกแบบขึ้นเพื่อทำหน้าที่เป็นแพลตฟอร์มการประเมินสุขภาพดินและวินิจฉัยสุขภาพพืชแบบฝังตัวบนขอบเครือข่าย (On-Device \& Edge Computing) ที่ทำงานได้แบบออฟไลน์ 100\% โดยผสมผสานการวิเคราะห์ข้อมูลภาพถ่ายภาคสนามความแม่นยำสูง (Agri-Vision 2.0 Precision ROI), การประมวลผลข้อมูลสเปกโทรสโกปีคลื่นแสงสะท้อน (Reflectance Spectroscopy), โครงข่ายเซนเซอร์เคมีดินและโทรมาตรฟาร์มอัจฉริยะ (Edge AIoT Telemetry) เพื่อให้คำแนะนำการจัดการดินและพืชที่แม่นยำและฉับไวต่อสถานการณ์จริงในแปลงเพาะปลูก

\academicsection{ความสำคัญของข้อมูลต่อการพัฒนาปัญญาประดิษฐ์}
\label{sec:data_importance}

ในกระบวนการเรียนรู้ของเครื่อง (Machine Learning) และการเรียนรู้เชิงลึก (Deep Learning) ประสิทธิภาพและความแม่นยำของโมเดลขึ้นอยู่กับคุณภาพของข้อมูลฝึกฝน (Training Data) เป็นปัจจัยชี้ขาด ความผิดพลาดที่พบบ่อยในการจัดฝึกอบรมเชิงปฏิบัติการทางปัญญาประดิษฐ์คือ การปล่อยให้ผู้เรียนจัดเก็บข้อมูลภาพเองอย่างกระจัดกระจาย ซึ่งนำไปสู่ปัญหาคุณภาพข้อมูลต่ำ ป้ายกำกับผิดพลาด และความไม่สมดุลของข้อมูล (Class Imbalance)

การจัดเตรียมชุดข้อมูลมาตรฐานที่มีโครงสร้างชัดเจน พร้อมทั้งจำแนกแหล่งที่มาของข้อมูลออกเป็นข้อมูลจริงที่ผ่านการรับรองทางวิชาการ (Real Benchmark Data) และข้อมูลจำลองเสมือนจริง (Synthetic Simulation Data) จึงเป็นขั้นตอนหัวใจสำคัญที่ช่วยให้ผู้เรียนสามารถทำความเข้าใจวงจรชีวิตการพัฒนา AI ได้อย่างถูกต้อง ตั้งแต่การเตรียมข้อมูล การฝึกโมเดล ไปจนถึงการประเมินผลความแม่นยำ

\academicsection{ขอบเขตและเป้าหมายของคู่มือปฏิบัติการ}
\label{sec:objectives}

คู่มือปฏิบัติการเล่มนี้จัดทำขึ้นโดยมีวัตถุประสงค์หลัก 4 ประการ ดังนี้

\begin{enumerate}
    \item เพื่ออธิบายสถาปัตยกรรมการจัดหมวดหมู่ข้อมูลทางการเกษตรแบบสองเสาหลัก (Two-Pillar Architecture) ได้แก่ โดเมนดิน (Soil Domain) และโดเมนพืช (Plant Domain)
    \item เพื่อนำเสนอแนวทางการแยกแยะข้อมูลจริงและข้อมูลจำลองในทุกโฟลเดอร์พืชและดินอย่างเป็นระบบ ตามหลักการวิจัยที่โปร่งใสและตรวจสอบได้
    \item เพื่อให้แนวปฏิบัติการฝึกฝนโมเดลจำแนกภาพด้วยสถาปัตยกรรม MobileNetV3 และการส่งออกเป็นมาตรฐานเปิด ONNX สำหรับใช้งานบน Edge AI และเว็บแอปพลิเคชัน
    \item เพื่อรวบรวมแคตตาล็อกฐานข้อมูลเปิดระดับโลกที่สำคัญด้านดินและพืช สำหรับใช้เป็นแหล่งอ้างอิงและเครื่องมือในการต่อยอดงานวิจัยระดับสูง
\end{enumerate}

ในบทถัดไปจะกล่าวถึงผังโครงสร้างสถาปัตยกรรมชุดข้อมูล และการบริหารจัดการไฟล์ภาพถ่ายจริงร่วมกับไฟล์ภาพจำลองในแต่ละโฟลเดอร์ย่อยอย่างละเอียด
